% Corpo comum dos testes: exercita capa, folhas, resumo, listas, siglas, ambientes e referências
\begin{document}
\pretextual
\imprimircapa
\imprimirfolhaderosto
\imprimirfolhadeaprovacao
\begin{resumo}
 \noindent Texto do resumo.

 \palavraschave{teste, classe, abntex2}
\end{resumo}
\begin{resumo}[Abstract]
\begin{otherlanguage*}{english}
   \noindent Abstract text.

   \keywords{test, class, abntex2}
 \end{otherlanguage*}
\end{resumo}
\listoffigures*
\cleardoublepage
\listoftables*
\cleardoublepage
\listadesiglas
\tableofcontents*
\textual
\chapter{Introdução}\label{cap:intro}
Texto com sigla (\ac{abnt}) e novamente \ac{abnt}; e \ac{api}. Veja a \autoref{fig:teste} e a
\autoref{tab:teste}.
\ifdefined\printbibliography Citação: \cite{teste}.\fi

\begin{figure}[htbp]
\begin{ilustracao}{\caption{Figura de teste}\label{fig:teste}}{\elaboracaopropria}
  \rule{4cm}{2cm}
\end{ilustracao}
\end{figure}

\begin{table}[htbp]
\begin{tabela}{\caption{Tabela de teste}\label{tab:teste}}{\fonte{Elaborada pelo autor.}}
  \begin{tabular}{lr}
    \toprule
    Item & Valor \\
    \midrule
    A & 1,0 \\
    B & 2,5 \\
    \bottomrule
  \end{tabular}
\end{tabela}
\end{table}

\ifdefined\printbibliography
\postextual
\printbibliography
\fi
\end{document}
